\section{AT Beweise der Inspiration}

Zu Beginn folgende Frage?\\
\betonung[p, b]{Glaubt ihr, dass dieses Buch Gottes Wort ist?} \\
Weil wenn nicht, dann bringen euch die nachfolgenden Beweise nichts. Diese stammen alle aus der Bibel.

Das ist wichtig, wenn ihr mit Leuten sprecht, die nicht glauben, dass die Bibel das Wort Gottes ist. Denen müsst ihr zuerst den Beweis erbringen, dass alles was in der Bibel steht, absolut unfehlbar ist. Nur kann das niemand beweisen. 

Wir können aufzeigen, dass es ein besonderes Buch ist. Dass dieses Buch bis jetzt existiert ist ein Wunder ist, dass es nachweislich erfüllte Voraussagen in der Bibel gibt.\\
Das alles können wir versuchen zu erklären und den Leuten nahe bringen, aber ob die Person auch wirklich an die Unfehlbarkeit der Bibel glaubt, dass können wir nicht erzwingen. Es gibt viele Christen die nur die Teile aus der Bibel nehmen, den ihnen angenehm sind. Gewöhnlich fängt das Streichen aus der Bibel schon auf der ersten Seite an. So werden meine Ausführungen auch nur den Leute als Beweis dienen, die an die Unfehlbarkeit der Bibel glauben.

Beweise das die Bibel Gotteswort ist, sind im Alten Testament vielfältig. Ich habe hier vier Punkte aufgelistet:
\begin{enumerate}
  \item Das Wort Gottes
  \item Direkte Reden Gottes
  \item Prophetische Reden Gottes
  \item Diktierte Texte von Gott.
\end{enumerate}

\subsection{Das Alte Testament wird bezeichnet als \enquote{die Worte Gottes}}
An etlichen Stellen im Alten Testemant wird geschrieben, dass dies das Wort Gottes ist. Es gibt ein Psalm mit 176 Versen, in dem alleine 24x geschrieben wird dass dies das Wort Gottes ist. Die Schrift wird in den 176 Versen, 175x erwähnt. Aber nicht nur dort, sondern auch in anderen Texten in der Bibel wird das bestätigt.

So zum Beispiel wird in diesen Bibeltexten konkret: \enquote{Gott sprach} genannt.\\
\begin{itemize}
  \item \biblezit{IIMos}{17:14}
  \item \biblezit{IIMos}{19:3.6-7}
  \item \biblezit{IKoen}{22:19}
  \item \biblezit{IIKoen}{20:16}
\end{itemize}

\subsection{Das Alte Testament enthält Worte Gottes in direkter Rede}
Vom direkten Reden von Gott lesen wir schon in den ersten Versen in der Bibel.
\begin{bibelquote}{SCHL}{IMos}{1:3}
  Und Gott sprach: \enquote{Es werde Licht! Und es wurde Licht.}
\end{bibelquote}
Auch lesen wir da von von Gottes Urteilen über Handlungen samt Ankündigung über entsprechende Folgen.
\begin{bibelquote}{SCHL}{IMos}{3:13-19}
  Da sprach Gott zu der Frau:...
\end{bibelquote}
Gott redet direkt mit Adam und Eva und zeigt ihnen auf, welche Auswirkungen ihre Tat hat.

Gott rief Abraham aus Ur und spricht auch mit Mose persönlich aus einem brennenden Dornbusch an.

Im ganzen Alten Testament lesen wir, dass Gott mit den Propheten gesprochen hat. Einige dieser Offenbarungen sind verbal,
\begin{bibelquote}{SCHL}{ISam}{3:21}
  Und der Herr erschien weiterhin in Silo; denn der Herr offenbarte sich dem Samuel in Silo durch das Wort des Herrn.
\end{bibelquote}
andere geschahen durch Visionen und Träume.
\begin{bibelquote}{SCHL}{IKoen}{3:5}
  In Gibeon erschien der Herr dem Salomo bei Nacht im Traum. Und Gott sprach: \enquote{Bitte, was ich dir geben soll!}
\end{bibelquote}
All das sind Aufzeichnungen göttlicher Rede.
\subsection{Das Alte Testament enthält prophetische Reden vonseiten Gottes}
Die Propheten wurden von Gott autorisiert für ihn zu sprechen. Und das, was sie sagten wurde als Gottes Reden angesehen. So sollte Mose zum Pharao gehen und ihm die Worte Gottes, mit den Worten \enquote{so spricht Gott} mitteilen.
\begin{bibelquote}{SCHL}{IIMos}{4:21}
  Und der Herr sprach zu Mose: \enquote{Wenn du wieder nach Ägypten kommst, so achte darauf, dass du vor dem Pharao all die Wunder tust, die ich in deine Hand gegeben habe. Ich aber will sein Herz verstocken, dass er das Volk nicht ziehen lassen wird.}
\end{bibelquote}
Diese Muster zieht sich durch das ganze Alte Testament. Wenn ein Prophet im Namen Gottes spricht, so ist die dabei verwendete typische Botenformel \enquote{So spricht der Herr}. Der Prophet kann sogar in der ersten Person für Gott reden und David erkannte an, dass Gott durch ihn sprach.
\begin{bibelquote}{SCHL}{IISam}{23:2}
  Der Geist des Herrn hat durch mich gerdet, und sein Wort war auf meiner Zunge.
\end{bibelquote}
Da auch andere, von Gott nicht berufene Personen, diese Formel benutzen, hat Gott selbst ein Kriterium genannt, anhand dem man Prüfen kann, ob der Überbringer der Botschaft wirklich Gottes Wort weiter gibt. So zu lesen in \biblezit{VMos}{13:1-6} und die Konsequenz war:
\begin{bibelquote}{SCHL}{VMos}{18:20}
  Doch der Prophet, der so vermessen ist, in meinem Namen zu reden, was ich ihm nicht zu reden geboten habe, oder der im Namen anderer Götter redet, jener Prophet soll sterben!
\end{bibelquote}
Mit diesem Kriterium kann man noch heute prüfen, ob die Worte ausgedacht sind oder ob diese wirklich von Gott kommen.
\subsection{Das Alte Testament enthält von Gott diktierte Reden}
Verschiedene Berichte in der Bibel werden von Gott diktiert. Fangen wir wieder bei Mose an. So lesen wir, dass Gott Mose die Anweisung gab seine Worte niederzuschreiben.
\begin{bibelquote}{SCHL}{IIMos}{34:27}
  Und der HERR sprach zu Mose: Schreibe dir diese Worte auf! Denn aufgrund dieser Worte habe ich mit dir und mit Israel einen Bund gemacht.
\end{bibelquote}
Auch am Ende seines Lebens wurde ihm befohlen, er solle die Gesetze niederschreiben, die ihm der Herr geboten hat. 
\begin{bibelquote}{SCHL}{VMos}{31:24-26}
  \hh{24}Als nun Mose damit fertig war, die Worte dieses Gesetzes vollständig in ein Buch zu schreiben, \hh{25}da gebot er den Leviten, welche die Bundeslade des HERRN trugen, und sprach: \hh{26}Nehmt das Buch dieses Gesetzes und legt es neben die Bundeslade des HERRN, eures Gottes, damit es dort ein Zeuge gegen dich sei.
\end{bibelquote}
Bei anderen Gelegenheiten wird er einfach von Gott angewiesen, aufzuzeichnen was geschehen war. 
\begin{bibelquote}{SCHL}{IIMos}{17:14}
  Da sprach der HERR zu Mose: Schreibe das zum Gedenken in ein Buch und präge es den Ohren Josuas ein: Ich will das Andenken Amaleks ganz und gar austilgen unter dem Himmel!
\end{bibelquote}
Als David seine Psalmen schrieb, wusste er, dass es Gott war, der durch ihn redete und trotzdem sind auch seine eigenen Gedanken, Worte und Empfindungen mit in die Texte eingeflossen.

Jeremia wurde angewiesen alle Worte aufzuschreiben, die Gott zu ihm redete.
\begin{bibelquote}{SCHL}{Jer}{30:1-4}
  \hh{1}Dies ist das Wort, das vom HERRN an Jeremia erging: \hh{2}So spricht der HERR, der Gott Israels: Schreibe dir alle Worte, die ich zu dir geredet habe, in ein Buch! \hh{3}Denn siehe, es kommen Tage, spricht der HERR, da ich das Geschick meines Volkes Israel und Juda wenden werde, spricht der HERR; und ich werde sie wieder in das Land zurückbringen, das ich ihren Vätern gegeben habe, und sie sollen es in Besitz nehmen. \hh{4}Das aber sind die Worte, die der HERR zu Israel und Juda gesprochen hat:
\end{bibelquote}

Jetzt könnt ihr selbest das Alte Testament durchforschen, wo Gott direkt redet. Aber nicht vergessen, auch wenn es  oft vorkommt, spricht Gott in der Geschichte eigentlich selten. Das Alte Testament dauert von der Schöpfung bis ca. 500 Jahre vor der Geburt Jesus. Es gibt ganze Jahrhunderte wo Gott kein Wort sagt.

Das nächste mal, es wird dann schon Mitte Dezember sein, werden wir im Neuen Testament nach Beweisen suchen, dass es Gottes Wort ist.




















